\ifdefined\tool\else
\documentclass[uplatex, 10pt, a4j, dvipdfmx]{jsarticle}

\usepackage{packages/abstract}
\usepackage{enumitem}
\usepackage[dvipdfmx]{graphicx}
\usepackage{url}
\usepackage[dvipdfmx, hidelinks]{hyperref}
\usepackage{pxjahyper}

\newcommand{\tool}{TOOL}
\newcommand{\AbstractStandalone}{}

% make abstract コマンドでの発表要旨のPDF生成に必要な記述項目
\setAbstractTitle{LLMを活用したシステム仕様書章立て生成のための\\プロンプトの作成と評価}
\setAbstractPresenter{西島尚紀}

\begin{document}
\AbstractHeader
\begin{AbstractBody}
\fi

ソフトウェア開発において、システム仕様書は開発の指針となる重要なドキュメントである。
特に、顧客の曖昧な要求を具体的な要件に落とし込む過程において、システム仕様書の品質はプロジェクトの成否を左右する。
しかし、システム仕様書に記述するべき内容は、組織や対象とするシステムの特性に依存する。
また、システム仕様書の質は、作成者の経験や知識によって左右される。

これらの課題を解決するために、LLM(大規模言語モデル)を用いたソフトウェア要求仕様書(本研究でのシステム仕様書に相当)生成の自動化を目的とした先行研究がある。
先行研究では、要件定義書や会議の議事録といった自然言語のドキュメントを入力として、情報を要約、抽出、整形し、あらかじめ用意した章立てに整形した情報を分類することで、ソフトウェア要求仕様書を自動生成する手法を提案している。
先行研究では、作成するソフトウェア要求仕様書の章立てはあらかじめ人間が用意することを前提としている。

しかし、前述したように、システム仕様書に記述するべき内容は対象とするシステムの特性に依存するため、適切な章立てを用意すること自体が、高度な知識と経験を要する困難な作業である。
また、システム仕様書の骨格となる章立てに欠陥があれば、後続の文章生成がいかに精巧であっても、システム仕様書としての品質は担保されない。
したがって、システム仕様書作成の自動化においては、用意された章立てへの情報の流し込みだけでなく、システムに応じた章立てそのものを生成するアプローチが必要となる。

そこで本研究では、「章立ての作成」に着目し、システム仕様書作成の対象となるシステムの特性に合わせた、章立ての自動生成を目的として、LLMを活用したシステム仕様書章立て生成のためのプロンプトを提案する。
システム仕様書章立て生成プロンプトは、「役割付与」、「ベースの章立て」、「開発するシステム情報」、「システム仕様書の定義」、「章立て出力指示」の5つの要素で構成する。

提案するプロンプトの有用性を確認するために、「BMP画像分割アプリ」と「クラウド型電子契約サービス」の2つのシステムを対象として、システム仕様書の章立てを生成した。
「BMP画像分割アプリ」は機能が少ない単純なシステムとして、「クラウド型電子契約サービス」は機能が多い複雑なシステムとして、実験対象とした。
また、作成した章立ての品質をシステム仕様書作成経験が5年以上ある専門家5名が評価者として、評価した。

評価項目は、「章立てに最低限含むべき項目を含んでいるか(網羅性)」、「章立ての順序は適切か(論理性)」、「章タイトルは適切か(命名の適切さ)」、
「章立て全体として理解しやすいか(可読性)」、「章立ての粒度は適切か(粒度の適切さ)」、「章の粒度は適切か(粒度の適切さ)」、
「想定読者に適した章立てか(適合性)」、「将来的な拡張に対応しやすい章立てか(拡張性)」、「章、節、項の包含関係は適切か(構造的整合性)」、
「章立てがシステム独自の章立てになっているか(独自性)」の9項目である。

実験の結果、提案プロンプトは、対象システムの複雑さに関わらず章立て品質のばらつきを抑制し、「品質の安定化」に寄与することを示した。
具体的には、提案プロンプトを用いない場合では、システムの特性により9項目の評価項目の合計点に3.8点の差が生じたのに対し、提案プロンプトを用いた場合は、差が2.3点であった。
これにより、単純なシステムにおける「検討項目の抜け漏れ」や、複雑なシステムにおける「専門知識不足による独自性の欠如」といった課題を、役割付与やベース章立ての提示によって効果的に補完し、一定水準の品質を確保可能であることを確認した。
また、更なる高品質な章立ての生成に向けては、プロンプトに含む指示内容の調整や、他のプロンプトエンジニアリング手法の導入が必要であると考察した。

\ifdefined\AbstractStandalone
\end{AbstractBody}
\end{document}
\fi
